\begin{claim}{C.26}{Proved}{Transitive centralizers on a closed hyperbolic surface}
Let $X$ be a closed connected orientable surface of genus $g\geq2$,
and give $A=\operatorname{Homeo}(X)$ the compact-open topology.
For every subgroup $L\leq A$ acting transitively on $X$, the centralizer
$C_A(L)$ is finite. Consequently, for an arbitrary strict action
$u:\Gamma\to A$ and a based loop $f_t$ in $C_{A^0}(u(\Gamma))$,
the loop is equivariantly null-homotopic whenever the subgroup generated
by its values acts transitively on $X$.
There also exists a closed torsion-free Polish subgroup
$G\leq\operatorname{Homeo}(D^2,\partial D^2)$ with $\pi_1G=\Z$;
extension by the identity embeds it closedly in every hyperbolic
$A^0$, where its displayed generator becomes null-homotopic. This
last subgroup is not asserted to be a centralizer.
\end{claim}
\begin{proof}
Put $K=\overline L$ and $N=C_A(K)$. Continuity of composition gives
$C_A(L)=C_A(K)$, and $K$ is a closed Polish subgroup of $A$ that
still acts transitively. Fix $x\in X$ and write $K_x$ for its stabilizer.
The open mapping principle for Polish actions gives a homeomorphism
\[
 K/K_x\longrightarrow X,\qquad kK_x\longmapsto kx.
                                                        \tag{C.26.1}
\]
Precisely, we use the ordinary-action case of
\cite[Theorem~3.5]{C_GomezPinedoUzcategui2017}: the acting group is
Polish, $X$ is nonmeagre Hausdorff, the action is continuous and
transitive, and the domain $K^x$ of the partial-action notation is all
of $K$, hence open. The source proves this open mapping statement;
no local section of (C.26.1) is being assumed.

Every $n\in N$ is determined by $n(x)$, since
$n(kx)=k n(x)$. Its possible values form exactly the closed subset
\[
 N_K(K_x)/K_x\ \subseteq\ K/K_x=X.                      \tag{C.26.2}
\]
Indeed, $n(x)=kx$ forces $K_x=K_{kx}=kK_xk^{-1}$.
Conversely a normalizing $k$ defines the homeomorphism
$gK_x\mapsto gkK_x$, commuting with $K$. The normalizer of the closed
subgroup $K_x$ is closed: both inclusions in
$kK_xk^{-1}=K_x$ are intersections of closed conditions. Its image
under (C.26.1) is closed because its inverse image under the quotient
map is that same normalizer. Thus (C.26.2) is compact.

Evaluation is a homeomorphism of $N$ onto (C.26.2). Only inverse
continuity needs proof. Fix $\varepsilon>0$. For every $y=g_yx$ choose
an identity neighborhood $W_y$ in $K$ such that each
$g_ywg_y^{-1}$, $w\in W_y$, moves every point of $X$ by less than
$\varepsilon/3$. The open sets $g_yW_yx$ cover $X$; take a finite
subcover with elements $g_1,\ldots,g_m$. When $n(x)$ is sufficiently
close to $n_0(x)$, each $g_i n(x)$ is within $\varepsilon/3$ of
$g_i n_0(x)$. At $z=g_iwx$ in the corresponding set, commutation gives
\[
 n(z)=g_iw n(x),\qquad n_0(z)=g_iw n_0(x).
\]
The two uniformly small movements by $g_iwg_i^{-1}$ and the middle
$\varepsilon/3$ bound imply $d(n(z),n_0(z))<\varepsilon$. This is
uniform in $z$ and proves inverse continuity in the compact-open
topology. In particular $N$ is compact.

An element of $N$ fixing one point fixes every point, by transitivity
and commutation. Every element of $A^0$ has Lefschetz number
$\chi(X)=2-2g\ne0$, hence has a fixed point. Therefore
\[
                       N\cap A^0=1.
\]
Projection to the discrete group $\pi_0A$ is thus a continuous
injection on $N$. Its image is compact and discrete, hence finite,
as asserted.

For the stated loop, put $L=\langle f_t:0\leq t\leq1\rangle$.
Every $u(\gamma)$ commutes with $L$, so the transitivity hypothesis
and the first assertion place $u(\Gamma)$ inside a finite group.
Write $H=u(\Gamma)$. The equality
$C_{A^0}(u(\Gamma))=C_{A^0}(H)$ is equality of subspaces, irrespective
of the cardinality of $\Gamma$. C.7 applied to the finite group $H$
contracts the loop in this exact centralizer. The assertion therefore
does not replace the operator group by a different action without a
comparison map.
\end{proof}

\paragraph{A proved obstruction to a closed-subgroup-only argument.}
There is a closed torsion-free Polish subgroup of
$\operatorname{Homeo}(D^2,\partial D^2)$ with fundamental group
$\Z$. We give the embeddings, because this example explains exactly
why closedness and absence of compact circles do not by themselves
settle C.12. It is not asserted to be a centralizer.

Let
\[
 E=\{f\in L^1([0,1],\R):f\text{ is integer-valued almost everywhere}\}.
\]
An $L^1$-convergent sequence has an almost-everywhere convergent
subsequence; hence $E$ is closed and Polish. For $0<s\leq1$ define
$D_sf(t)=f(t/s)$ on $[0,s]$ and zero on $(s,1]$, and put $D_0f=0$.
These maps remain in $E$ and satisfy $\|D_sf\|_1=s\|f\|_1$.
Strong continuity of dilation follows first for step functions and
then by $L^1$ approximation. Thus $(s,f)\mapsto D_sf$ contracts $E$
and each of its norm balls at zero. In particular $E$ is contractible
and locally contractible.

The constant integer functions form a closed discrete subgroup $\Z$.
The quotient $G=E/\Z$ is Polish, and $E\to G$ is a covering: translates
of a norm ball of radius less than $1/3$ by distinct integers are
disjoint. It is the universal cover, and
\[
 \pi_1(G)=\Z,\qquad
 \gamma(s)=[\mathbf1_{[0,s]}]\quad(0\leq s\leq1)          \tag{C.26.3}
\]
is a generator, because its lift joins zero to the constant function
$1$. Nevertheless $G$ is torsion-free: $n[f]=0$ says that $nf$ is
a constant integer, so the integer-valued $f$ is constant and $[f]=0$.

There is a closed isometric group embedding
\[
 G\longrightarrow V:=L^1([0,1],\R)/\R,
       \qquad[f]_{\Z}\longmapsto[f]_{\R},                 \tag{C.26.4}
\]
where the second denominator consists of all constant real functions.
For integer-valued $h$, the convex function $c\mapsto\int|h-c|$ is
linear between consecutive integers, tends to infinity at both ends,
and attains a minimum at an integer. Hence
\[
 \inf_{c\in\R}\|h-c\|_1=\inf_{n\in\Z}\|h-n\|_1,
\]
proving the isometry. The quotient metric on $G$ is complete: choose
a Cauchy subsequence with successive quotient distances below $2^{-j}$
and adjust each representative by an integer so that the successive
$L^1$ distances have that bound. The representatives converge in $E$,
and the entire original Cauchy sequence then converges in $G$.
This proves closedness in the Banach space $V$.

We spell out a closed linear isometric embedding
$T:V\hookrightarrow C_0([0,1],\R)$, where the subscript here means
vanishing at both endpoints. The weak-star unit ball $B$ of $V^*$ is
compact metrizable, and $v\mapsto(\ell\mapsto\ell(v))$ embeds $V$
isometrically into $C(B)$. A continuous surjection from the Cantor
set to $B$ exists by nested finite covers of decreasing diameters.
Pullback followed by linear interpolation across each complementary
interval of the Cantor set embeds $C(B)$ linearly and isometrically
in $C([0,1])$. Place this interval in the middle third and interpolate
to zero at the two outer endpoints. The resulting map is the required
$T$; its image is closed because $V$ is complete.

For $F\in C_0([0,1])$, the formula
\[
 \Phi_F(x,y)=
   \bigl(\tanh(\operatorname{artanh}(x)+F(y)),y\bigr)
          \qquad(|x|<1)
                                                        \tag{C.26.5}
\]
on $[-1,1]\times[0,1]$, extended by the identity on the vertical
sides, fixes the entire boundary and satisfies
$\Phi_F\Phi_Q=\Phi_{F+Q}$. It is a topological-group embedding:
$\tanh$ is Lipschitz, giving continuity, and evaluation at $(0,y)$
recovers $\tanh(F(y))$, from which $F$ is recovered uniformly near
any given profile. Its image is closed. A homeomorphism limit fixes
the boundary and preserves the horizontal leaves; the image of the
compact segment $\{0\}\times[0,1]$ stays a positive distance from the
vertical sides. Applying $\operatorname{artanh}$ therefore recovers
uniformly bounded, uniformly convergent profiles.

Composing (C.26.4), $T$, and (C.26.5) embeds $G$ as a closed subgroup
of $\operatorname{Homeo}(D^2,\partial D^2)$. Extension by the identity
outside a coordinate disc is a closed topological-group embedding
into $A^0$ for any hyperbolic surface: restriction recovers the maps,
and being the identity off the disc is a closed condition. Its loop
(C.26.3) is ambient-null but nontrivial in $G$.
The precise enlargement killing it is already visible: if
$F_s=T([\mathbf1_{[0,s]}]_{\R})$, then $F_0=F_1=0$ and
$(r,s)\mapsto\Phi_{rF_s}$ is a based contraction in the larger real
shear group. No operator group cutting out a centralizer that excludes
this contraction has been constructed here. Thus this example refutes
the auxiliary assertion of fundamental-group injectivity for all
closed subgroups, while C.12 retains its stated centralizer question.
